% =================================================================
% 1. 文档类定义
% =================================================================
\documentclass[acmlarge,nonacm]{acmart}

% =================================================================
% 2. 核心与基础包
% =================================================================
%% BibTeX 命令保护 (好习惯)
\AtBeginDocument{%
  \providecommand\BibTeX{{Bib\TeX}}%
}
\usepackage{graphicx}  % acmart已加载, 但显式加载无害
\usepackage{xcolor}    % acmart已加载, 但显式加载无害
\usepackage{textcomp}

% =================================================================
% 3. 数学相关包 (修正后的加载顺序)
% =================================================================
\usepackage{mathtools} % 自动加载 amsmath

% [推荐] 只加载 stmaryrd 中需要的特定符号以避免冲突
\usepackage[only,llbracket,rrbracket]{stmaryrd} 

\usepackage{siunitx}

% =================================================================
% 4. 表格与格式化包
% =================================================================
\usepackage{booktabs}
\usepackage{adjustbox}
\usepackage{threeparttable}
\usepackage{multirow}
\usepackage{makecell}
\usepackage{array}
\usepackage{ragged2e}
\usepackage{tabularx}
\newcolumntype{L}[1]{>{\RaggedRight\arraybackslash}p{#1}}

% =================================================================
% 5. 列表、代码与算法包
% =================================================================
\usepackage{listings}
\usepackage[inline]{enumitem}
\usepackage{algorithm}
\usepackage{algorithmicx}
\usepackage{algpseudocode}

% =================================================================
% 6. 定理环境 (acmart 已加载 amsthm, 直接定义即可)
% =================================================================
\theoremstyle{acmplain}
\newtheorem{theorem}{Theorem}
\newtheorem{lemma}[theorem]{Lemma}
\newtheorem{corollary}[theorem]{Corollary}
\newtheorem{proposition}[theorem]{Proposition}

\theoremstyle{acmdefinition}
\newtheorem{definition}{Definition}[section]
\newtheorem{example}{Example}[section]
\newtheorem{remark}{Remark}[section]

% =================================================================
% 7. 布局与自定义命令
% =================================================================
\usepackage[section]{placeins} % 控制浮动体位置
\graphicspath{{figures/}}      % 设置图片路径 (好习惯)

%% 自定义数学算符 (好习惯)
\DeclareMathOperator*{\argmax}{arg\,max}
\DeclareMathOperator*{\argmin}{arg\,min}

%% 代码列表设置
\lstset{
  basicstyle=\footnotesize\ttfamily,
  numbers=none,
  numberstyle=\tiny,
  frame=none,
  breaklines=true,
  captionpos=b,
  showstringspaces=false,
  commentstyle=\color{gray},
  keywordstyle=\color{blue},
  stringstyle=\color{red}
}

\begin{document}

%% -----------------------------------------------------------------
%% TITLE & AUTHORS
%% -----------------------------------------------------------------

\title{ATLAS: A Layered Constraint-Guided Framework for Structured Artifact Generation in LLM-Assisted MDE}


\author{Tong Ma}
\orcid{0009-0000-6481-0908}
\affiliation{%
  \institution{University of Science and Technology of China}
  \city{Hefei}
  \country{China}
}
\affiliation{%
  \institution{Hefei Institutes of Physical Science, Chinese Academy of Sciences}
  \city{Hefei}
  \country{China}
}
\email{matong@mail.ustc.edu.cn}

\author{Hui Lai}
\affiliation{%
  \institution{University of Science and Technology of China}
  \city{Hefei}
  \country{China}
}

\author{Hui Wang}
\affiliation{%
  \institution{Anhui University}
  \city{Hefei}
  \country{China}
}

\author{Zhenhu Tian}
\affiliation{%
  \institution{Anhui University}
  \city{Hefei}
  \country{China}
}

\author{Chaochao Li}
\affiliation{%
  \institution{University of Science and Technology of China}
  \city{Hefei}
  \country{China}
}
\authornote{Corresponding author.}

\author{Fengjie Xu}
\orcid{0009-0004-5185-0619}
\affiliation{%
  \institution{Hefei Institutes of Physical Science, Chinese Academy of Sciences}
  \city{Hefei}
  \country{China}
}
\authornote{Corresponding author.}

\author{Ling Fang}
\orcid{0000-0003-3376-3197}
\affiliation{%
  \institution{Hefei Institutes of Physical Science, Chinese Academy of Sciences}
  \city{Hefei}
  \country{China}
}
\authornote{Corresponding author.}
\email{fangl@hfcas.ac}

%% Short author list for headers
\renewcommand{\shortauthors}{Ma et al.}



%% -----------------------------------------------------------------
%% ABSTRACT
%% -----------------------------------------------------------------

\begin{abstract}
ATLAS is a constraint-guided generation framework for structured engineering artifacts whose outputs must satisfy explicit schemas, domain rules, and audit requirements. Rather than treating a large language model as a standalone generator, ATLAS places generation inside a model-driven workflow that separates domain representation, constraint compilation, and post-generation validation.

ATLAS combines three components. A metamodel-integration stage builds a typed representation of domain entities and relations; in this study, it operates over authoritative AUTOSAR meta-model assets. An Integrated Constraint Model (ICM) compiles heterogeneous requirements into two operational layers: generation-time structural constraints and post-generation semantic/logical obligations. Constraint-Guided, Validation-Backed Generation (CVG) then combines Layer~1 constrained decoding, Layer~2 backend validation, and audit-guided repair. In the AUTOSAR instantiation, these Layer~2 obligations are realized through SHACL/SMT-style checks, illustrating how the same ICM can be connected to domain-specific validation backends.

We evaluate ATLAS on AUTOSAR artifact generation at both single-file and multi-file scales. In the evaluated AUTOSAR setting, ATLAS consistently produces schema-valid single-file outputs and preserves perfect file completeness and XSD validity at multi-file scale, while SHACL/SMT checks and result analysis continue to expose residual system-level defects. The empirical picture is therefore one of bounded automation: ATLAS secures structural validity and turns higher-level failures into explicit, diagnosable objects within the generation workflow.
\end{abstract}



%% -----------------------------------------------------------------
%% KEYWORDS
%% -----------------------------------------------------------------

\keywords{
model-driven engineering;
metamodel integration;
constraint-guided generation;
large language models;
constraint validation;
auditable artifacts
}

\maketitle



\section{Introduction}
\label{sec:introduction}

Large Language Models (LLMs) have made it increasingly practical to generate code, configuration files, and other structured artifacts directly from natural-language specifications. This capability is attractive in safety- and compliance-critical settings, where engineers routinely translate requirements into artifacts that must conform to domain schemas, preserve traceability, and survive external review. Yet this is also where unconstrained generation is least acceptable: a syntactically fluent output is not enough if it violates the governing meta-model, breaks cross-reference integrity, or leaves no inspectable explanation of why it should be trusted.

Model-Driven Engineering (MDE) addresses part of this problem well. Mature MDE pipelines provide explicit meta-models, typed transformations, and validation mechanisms that improve repeatability and make structural correctness checkable~\cite{omg2003mda,kahaniSurveyClassificationModel2019,delavaraModelbasedSpecificationSafety2016}. In automotive software and related schema-rich engineering settings, these assets are especially valuable because the artifacts being produced are already embedded in a broader assurance process. However, conventional MDE pipelines are typically engineered around stable schemas, carefully maintained transformations, and domain-specific toolchains. They are strong on determinism, but comparatively weak at handling underspecified natural-language inputs and rapidly evolving artifact requests.

Recent work suggests that LLMs and MDE can complement each other rather than compete. LLMs are effective at mapping informal descriptions into candidate structures, while MDE contributes the typed representations and validation boundaries needed to keep those structures meaningful~\cite{Petrovic_2024,Alaoui2025,Patil2025}. At the same time, constrained decoding, retrieval grounding, and validator-in-the-loop repair have shown that generation becomes substantially more reliable when the model is forced to respect external structure during or after decoding~\cite{scholakPICARDParsingIncrementally2021,poesiaSynchromeshReliableCode2021,gengGrammarConstrainedDecodingStructured2024,zhangKnowGPTKnowledgeGraph2024,tangCodeRepairLLMs2024,shinn2023reflexion,chen2024selfdebug}. In AUTOSAR-style settings, what is still missing is a framework that turns these ingredients into a single engineering pipeline: one that clearly separates the domain model from the generated artifact, distinguishes structural from semantic constraints, and records evidence that can be inspected when automation falls short.

This paper presents ATLAS, a layered framework for evidence-oriented artifact generation. ATLAS is organized around three ideas. First, it integrates authoritative schemas into a typed domain metamodel that serves as the semantic backbone for downstream generation, while leaving room for future text-assisted schema induction when machine-readable schemas are unavailable. Second, it compiles heterogeneous requirements into an \emph{Integrated Constraint Model (ICM)} with two operational layers: structural constraints that can be enforced during decoding, and semantic/logical obligations that are evaluated after generation. Third, it executes generation through \emph{Constraint-Guided, Validation-Backed Generation (CVG)}, which combines prefix-safe structural masking, post-generation validation, and audit-guided repair.

This separation reflects the structure of the engineering problem itself. Some requirements are local enough to be enforced during decoding, whereas others depend on completed artifacts, cross-file context, or domain-specific analysis and therefore emerge only at validation time. ATLAS is designed around that asymmetry: it uses Layer~1 to preserve structural admissibility during generation, and uses Layer~2 to connect the generated artifact to whichever semantic and logical checks are meaningful for the target domain. In domains with mature validation toolchains, this means that ATLAS can sit in front of existing validators; in domains with less complete infrastructure, the same boundary provides a place to introduce formal checks incrementally.

The paper makes three contributions:
\begin{enumerate}[leftmargin=*,itemsep=0.4ex]
    \item ATLAS provides a layered architecture for schema-guided artifact generation. We define it as a separation of metamodel integration, requirement compilation (ICM), and constraint-aware execution (CVG), so that schema acquisition, generation, validation, and repair can be reasoned about as distinct steps.
    \item ATLAS implements constraint-aware generation with explicit evidence flow. It enforces structural constraints during decoding, applies semantic and logical checks after generation through Layer~2 validation backends, and records validation outcomes and repair traces to support inspection and reuse.
    \item We conduct an empirical study centered on AUTOSAR. The evaluation covers both single-file and multi-file generation, where explicit schemas, layered constraints, and industrial validation tooling make the benefits and limits of the approach observable.
\end{enumerate}

The rest of the paper is organized as follows. Section~\ref{sec:related-work} positions ATLAS with respect to MDE, constrained generation, and validation-oriented LLM workflows. Section~\ref{sec:methodology} presents the metamodel integration, ICM, and CVG framework. Section~\ref{sec:evaluation} reports the AUTOSAR-centered empirical evaluation, and Section~\ref{sec:conclusion} discusses implications, limitations, and future work.


\section{Related Work}
\label{sec:related-work}

\subsection{MDE for structured artifact engineering in schema-rich settings}
Model-Driven Engineering provides the classical foundation for deriving structured artifacts from explicit domain models through model-to-model and model-to-text transformations~\cite{omg2003mda,kahaniSurveyClassificationModel2019,DualStageFramework}. In schema-rich engineering settings, this tradition is especially relevant because meta-models, schemas, and validation rules are already part of the engineering process. Prior work in automotive and other safety-critical domains shows the practical value of authoritative schemas, traceable transformations, and tool-supported conformance checking~\cite{delavaraModelbasedSpecificationSafety2016,el-gnainyAIEnhancedAUTOSARConfiguration2024a,nairExtendedSystematicLiterature2014}. ATLAS builds on this line of work, but shifts the focal problem: rather than assuming that structured source models already exist and only need deterministic transformation, we study how natural-language specifications can be converted into structured artifacts \emph{under} domain constraints.

\subsection{LLMs for modeling, schema-guided generation, and automotive software tasks}
A growing body of work explores how LLMs can assist modeling activities, including model instance creation, DSL-oriented modeling, and software generation in automotive contexts~\cite{Petrovic_2024,Alaoui2025,Patil2025,DiRocco2025LLMMDE}. Recent automotive studies have pushed this line further by combining LLMs with event-chain models, retrieval from vehicle signal catalogs, and verification-oriented software workflows~\cite{Kirchner2025AutoCode,Petrovic2025SurveyAutomotiveGenAI,Petrovic2025EventChainADAS}. Taken together, these papers strengthen the case that LLMs are useful when the input is underspecified or expressed in natural language, but they also reinforce a common limitation: raw prompting remains fragile once outputs must obey nontrivial industrial structure.

In parallel, constrained decoding research has demonstrated that grammars, automata, or schema-aware decoders can substantially improve structural correctness by restricting the model to admissible continuations during generation~\cite{scholakPICARDParsingIncrementally2021,poesiaSynchromeshReliableCode2021,gengGrammarConstrainedDecodingStructured2024,Dong2024XGrammar,Park2025,Nguyen2026UnifiedDecoding}. Retrieval-augmented prompting and knowledge-graph grounding further reduce hallucination by injecting domain facts or schema fragments into the context window~\cite{zhangKnowGPTKnowledgeGraph2024}. ATLAS differs from these threads by combining schema construction, layered constraint compilation, constrained execution, and reusable evidence packaging in a single end-to-end pipeline aimed at structured engineering artifacts rather than general text generation. In the present paper, that pipeline is evaluated only in the schema-rich AUTOSAR setting.

\subsection{Validation, repair, and evidence-oriented generation}
Post-generation validation and iterative repair are now common strategies for making LLM outputs more dependable. In software engineering and formal-methods-oriented settings, validators such as type checkers, graph constraints, and solver-backed analyses can identify defects that the model then repairs in a subsequent round~\cite{tangCodeRepairLLMs2024,shinn2023reflexion,chen2024selfdebug}. This pattern is compatible with longstanding ideas from formal verification and evidence-oriented engineering, where generated outputs are accompanied by re-checkable validation records rather than trusted solely because of the generator that produced them. ATLAS adopts this evidence-oriented perspective, but narrows its claim carefully: instead of presenting validation results as a full substitute for domain certification, we use them as explicit, inspectable support for downstream review. The resulting position is between rigid deterministic transformation and unconstrained LLM generation: structure is enforced as early as possible, richer semantic conditions are checked after generation, and repair is driven by validator feedback instead of free-form re-prompting alone.

\subsection{Positioning of ATLAS}
The closest prior efforts typically emphasize one part of the pipeline at a time: meta-model construction, constrained decoding, retrieval grounding, or repair. Table~\ref{tab:related-positioning} summarizes the papers most relevant to ATLAS at the level of \emph{pipeline design}, rather than only at the level of individual techniques.

\begin{table*}[!htbp]
\centering
\small
\caption{Positioning ATLAS against closely related LLM$\times$MDE or structure-aware generation workflows.}
\label{tab:related-positioning}
\begin{tabular}{@{}p{1.5cm}p{2.3cm}p{2.9cm}p{2.9cm}p{4.5cm}@{}}
\toprule
\textbf{Work} & \textbf{Primary substrate} & \textbf{When structure is enforced} & \textbf{Evidence / repair stance} & \textbf{Main distinction from ATLAS}\\
\midrule
Petrovi\'{c} et al.~\cite{Petrovic_2024} & Automotive meta-model instances & Mainly after generation via OCL consistency checks & Validation is present, but evidence is not packaged as a reusable audit object & Focuses on model-instance creation; does not separate metamodel integration / ICM / evidence as first-class layers\\
Kirchner and Knoll~\cite{Kirchner2025AutoCode} & Automotive software workflow with testing / verification loop & During downstream code and verification workflow & Verification-guided iteration over generated code & Targets code-generation workflow rather than constraint-compiled artifact generation\\
Petrovi\'{c} et al.~\cite{Petrovic2025EventChainADAS} & Event-chain + retrieved signal catalogs for ADAS & Through retrieved signals and event-chain validation & Validation is tied to the event-chain/code loop & Strong automotive relevance, but centered on event-driven code synthesis rather than general layered artifact validation\\
Al-Azzoni et al.~\cite{AlAzzoni2026SmartContractTest} & RBAC/MDE models for smart-contract testing & Model-guided test generation with downstream execution checks & Emphasizes executable test generation & Shows LLM$\times$MDE transferability, but not a two-layer constraint system for auditable structured artifacts\\
\textbf{ATLAS} & Integrated domain metamodel + compiled ICM + validation evidence & Both during decoding (L1) and after generation (L2) & Audit trail, validator outputs, and repair traces are explicit outputs & Integrates schema acquisition, layered enforcement, and inspectable evidence in one regulated-generation pipeline\\
\bottomrule
\end{tabular}
\end{table*}

ATLAS contributes by integrating these pieces around a common typed representation and by making the separation between \emph{domain model}, \emph{generated artifact}, and \emph{validation evidence} explicit. This separation is central to the framework design: the generated output is an artifact conformant to the domain metamodel, not the metamodel itself; structural validity and semantic acceptability are related but different properties; and human experts remain responsible for residual decisions that are not fully captured by the compiled constraints.



\section{Methodology: Metamodel Integration, ICM, and CVG}
\label{sec:methodology}

\subsection{Overview}

\begin{figure}[!htbp]
  \centering
  \includegraphics[width=0.85\linewidth,
                   trim=0.5cm 0.5cm 0.5cm 0.5cm,clip]{prism-overview-arch.pdf}
    \vspace{-1\baselineskip}
  \caption[Overview of the ATLAS framework.]{ATLAS links domain requirements to validation-backed generation. A metamodel-integration stage turns existing meta-model assets, and where necessary natural-language requirements, into an explicit domain metamodel and an Integrated Constraint Model (ICM). Generation is guided by structural constraints (Layer-1, enforced during decoding) and then checked by semantic / logical validators (Layer-2) after decoding. The validation signals, together with the decode-time audit record, form evidence that can be used to automatically repair the artifact in a closed loop.}
  \Description{Block diagram of the ATLAS pipeline. Inputs (NL requirements
  and meta-model specs) are integrated into an explicit domain metamodel and a constraint
  model. L1 automaton-guided decoding enforces structural constraints during
  generation, L2 semantic and logic validators check the final artifact, and the audit
  plus validation results feed an audit-guided repair loop.}
  \label{fig:prism-overview-arch}
  \vspace{-1\baselineskip}
\end{figure}

ATLAS is structured around a clear separation of roles among domain representation, artifact construction, and evidence production. The \emph{domain metamodel} is the domain-level representation that explains which entity types, attributes, and relations exist. The \emph{generated artifact} is an instance-level output that must conform to that representation. The \emph{evidence bundle} is a separate object that records how the artifact was generated and which checks it passed. This separation matters both technically and conceptually: it clarifies what is being constructed at each stage and prevents the framework from overstating what the system proves.

The pipeline has three components.

\paragraph{Metamodel integration.} This stage builds a provenance-aware, typed representation of the target domain. In domains with authoritative schemas, it is obtained by deterministic transformation; the architecture also leaves room for a text-assisted induction path when such schemas are unavailable. The resulting domain metamodel is the reference structure used for retrieval, validation anchoring, and human review.

\paragraph{Integrated Constraint Model (ICM).} The ICM is the persistent repository of constraints associated with the domain metamodel. It records where each constraint came from, which metamodel entities it is anchored to, how it is compiled, and whether it is intended for generation-time enforcement or post-generation validation. The key design choice is that LLM-extracted rules are never admitted directly: they must first be aligned to the metamodel and pass compatibility checks.

\paragraph{Constraint-guided, validation-backed generation (CVG).} CVG uses the ICM in two stages. Layer~1 (L1) constraints restrict decoding so that the draft remains structurally admissible. Layer~2 (L2) validators then check semantic and logical properties that require the completed artifact. The outputs of these validators, together with the decode-time audit trail, form an evidence bundle that supports targeted repair and downstream inspection.


\subsection{Authoritative Metamodel Integration and Extension Path}
\label{sec:metamodel-integration}

\paragraph{Problem setting.}
ATLAS must operate in two different starting conditions. In domains such as AUTOSAR, the engineering ecosystem already provides explicit artifacts such as XSD, XMI, or Ecore, and the core task is normalization. In domains without a canonical machine-readable schema, the system must first recover an operational domain vocabulary from prose. ATLAS supports both cases through a common metamodel representation, but the two construction paths should not be conflated: Path~S1 is authoritative and deterministic, whereas Path~S2 is an assisted induction process that requires stronger admission control.

\begin{definition}[Domain Metamodel as Typed Graph]
A domain metamodel in ATLAS is represented as a typed directed graph
$M = (V,E,T,\Phi)$, where $V$ is the set of domain entities, $E$ is the set of typed relations, $T$ assigns node and edge kinds, and $\Phi$ stores structural facts such as cardinalities, containment, typing, and mandatory references.
\end{definition}

This graph representation is not presented as a novel formalism. Its role is to give the later stages a clear interface: metamodel integration produces a typed graph, generation consumes subgraphs of that graph as context, and validation maps failures back to graph entities and source spans.

\begin{figure}[!htbp]
  \centering
  \includegraphics[width=0.6\linewidth,
    trim=0.5cm 0.5cm 0.5cm 0.5cm, clip]{umm-extraction.pdf}
  \caption[Dual-path metamodel integration.]{Two ways to build the domain metamodel used by ATLAS. Path S1 handles structured sources: when a machine-readable meta-model (for example AUTOSAR, AADL, OPC UA) already exists, ATLAS ingests and normalizes it. Path S2 handles unstructured sources: when only natural-language requirements are available, ATLAS uses an LLM to induce domain entities and relations directly from text. In this paper, S2 is presented as a design capability rather than an empirically evaluated deployment path.}
  \Description{Two-path process for building the domain metamodel used by ATLAS. 
  Path S1 applies when a machine-readable meta-model such as AUTOSAR XSD/XMI, 
  AADL, or OPC UA Ecore already exists: ATLAS ingests and normalizes that model 
  using model-to-model transformations. Path S2 applies when no authoritative 
  schema exists and requirements only appear in natural-language documents: 
  ATLAS uses an LLM-based extractor, guided by domain context, to induce 
  entities and relations. In this paper, S2 is retained as part of the 
  architecture but is not itself evaluated empirically.}
  \label{fig:umm-extraction}
  \vspace{-0.5\baselineskip}
\end{figure}

\paragraph{Path S1: transformation from authoritative schemas.}
When the domain already provides an explicit meta-model, ATLAS follows a deterministic transformation path adapted from our earlier S$^2$D$^2$ pipeline~\cite{DualStageFramework}. In our AUTOSAR instantiation, this path consists of three steps:
\begin{enumerate}[leftmargin=*,nosep]
  \item Structural extraction parses XMI/XSD artifacts to recover entities, inheritance relations, attributes, cardinalities, and reference structure.
  \item Normalization and fusion merge overlapping views into a single typed representation; where multiple schema sources describe the same concept, source provenance is preserved instead of being silently collapsed.
  \item Behavioral annotation carries forward operational cues that matter downstream, such as reference patterns, containment conventions, and domain-specific role semantics.
\end{enumerate}
Path~S1 therefore produces a \emph{normalized domain metamodel representation}. The later generation task creates artifacts that conform to the domain metamodel; it does not create the metamodel itself.

\paragraph{Path S2: cautious induction from text.}
When no authoritative schema exists, ATLAS uses an LLM to propose candidate entities, relations, and lightweight obligations from natural-language specifications. These proposals are not treated as authoritative by default. Instead, they are normalized into the same typed-graph format as Path~S1 and attached to provenance metadata (document, paragraph, confidence, extraction prompt version). The remainder of the pipeline therefore treats all text-derived structure as \emph{candidates} until grounded and checked. Because Path~S2 is not evaluated in the present AUTOSAR study, we retain it only to delineate the architectural scope of ATLAS and do not draw empirical conclusions from it here.

\begin{definition}[Partial Alignment]
\label{def:partial-align}
Let $U$ be the set of metamodel entities and relation signatures. For an extracted candidate $c$, the alignment map
\[
\alpha(c) \in U \cup \{\bot\}
\]
returns a metamodel anchor when the candidate can be mapped to an existing entity or relation with sufficient lexical and type compatibility; otherwise it returns $\bot$. Candidates with $\alpha(c)=\bot$ are not admitted into the persistent ICM and are instead quarantined for review.
\end{definition}

\paragraph{How the metamodel integration stage is used downstream.}
Once constructed, the domain metamodel serves three roles. It is serialized into retrieval snippets for prompting. It provides the anchors used by Channel~2 constraint extraction. It also drives the typed editing view used during human repair. This is why we keep the representation explicit in the paper: it is not merely a background artifact but the interface between deterministic engineering assets and LLM-facing generation.

\subsection{Constraint Extraction and Integrated Constraint Model (ICM)}
\label{sec:icm}

\paragraph{Constraint families and layer boundary.}
ATLAS organizes obligations into three operational families:

\begin{definition}[Operational Constraint Partition]
\label{def:constraint-space}
Given a domain metamodel, the admitted constraint repository is partitioned into
\[
\mathcal{R}
=
(\mathcal{R}_{\mathrm{struct}},
 \mathcal{R}_{\mathrm{sem}},
 \mathcal{R}_{\mathrm{log}}).
\]
$\mathcal{R}_{\mathrm{struct}}$ contains prefix-checkable structural obligations such as containment, field presence, ordering, and local typing patterns; $\mathcal{R}_{\mathrm{sem}}$ contains graph-level or reference-level obligations that require the completed artifact; and $\mathcal{R}_{\mathrm{log}}$ contains numeric, temporal, or solver-backed domain conditions. This partition is operational rather than ontological: one domain rule may contribute a structural fragment to L1 and a richer semantic or logical fragment to L2. The partition classifies obligations, not validator brands: the post-generation backend may be a formal encoding such as SHACL/SMT, a rule engine, a domain-specific checker, or an existing commercial configuration tool.
\end{definition}

This partition separates decode-time structural enforcement from post-generation semantic and logical validation. Prefix-checkable structure is enforced during decoding, while global constraints are checked only after artifact completion.

\paragraph{Channel 1 and Channel 2.}
Channel~1 extracts constraints from structured engineering assets and authoritative schemas. It is deterministic and supplies the initial structural backbone of the ICM. Channel~2 extracts candidate semantic and logic constraints from natural-language documents using an LLM, but every candidate must pass an admission workflow before it becomes persistent.

\begin{figure}[!htbp]
  \centering
  \includegraphics[width=0.7\linewidth,
    trim=0.5cm 1cm 0.6cm 0.5cm, clip]{doc-info-extraction.pdf}
  \caption[Channel 2 extraction with semantic validation.]{LLM-based extraction from natural-language specifications. The extractor proposes structured constraint candidates and aligns each candidate to entities in the domain metamodel. Only candidates that align cleanly and pass compatibility checks are committed to the Integrated Constraint Model (ICM); the rest are flagged for human review instead of being turned directly into validators.}
  \Description{Illustration of Channel 2 extraction with semantic validation.
  Specification text is processed by an LLM-based extractor, which proposes structured JSON-like constraint candidates and links each candidate to entities in the domain metamodel.
  Only candidates that can be anchored via alpha and that pass semantic compatibility checks are inserted into the Integrated Constraint Model (ICM).
  Candidates that fail these checks are not compiled directly into L2 validators but are instead flagged for review.}
  \label{fig:doc-info-extraction}
\end{figure}

\begin{definition}[Semantic Compatibility]
\label{def:semantic-compat}
A candidate constraint $c$ is semantically compatible with the current domain metamodel if: 
(i) its operands can be typed using the anchor returned by $\alpha(c)$ and the surrounding metamodel neighborhood;
(ii) it does not contradict already admitted structural facts such as cardinality, containment, or declared reference kinds; and
(iii) it can be compiled into the intended backend representation (e.g., SHACL, SMT, or a domain-validator interface) without leaving unresolved symbols.
\end{definition}

This definition operationalizes semantic compatibility through typed anchoring, contradiction checks, and successful backend compilation.

\begin{definition}[Integrated Constraint Model]
\label{def:icm}
The ICM is a tuple
\[
\mathcal{I}=(\mathcal{M},\mathcal{R},\mathcal{P},\Gamma),
\]
where $\mathcal{M}$ is the domain metamodel, $\mathcal{R}$ is the admitted constraint repository partitioned into $(\mathcal{R}_{\mathrm{struct}},\mathcal{R}_{\mathrm{sem}},\mathcal{R}_{\mathrm{log}})$, $\mathcal{P}$ stores provenance and compilation metadata for each constraint, and $\Gamma$ is a dependency order used to prioritize validation and repair.
\end{definition}

This definition specifies the interfaces that the ICM must support downstream: storage, provenance, compilation, and repair ordering.

\paragraph{ICM construction.}
Algorithm~\ref{alg_3.2:icm-construction} summarizes how the repository is populated. The key safeguard is that Channel~2 output either becomes a typed, compilable, provenance-carrying constraint or it stays out of the persistent repository.

\begin{algorithm}[t]
\caption{ICM Construction via Dual-Channel Extraction}
\label{alg_3.2:icm-construction}
\begin{algorithmic}[1]
\small
\Require Structured sources $S$, specification documents $D$, domain metamodel $\mathcal{M}$, constraint schema $\mathcal{S}_c$
\Ensure Persistent ICM $\mathcal{I}$

\Function{BuildICM}{$S,D,\mathcal{M},\mathcal{S}_c$}
  \State initialize empty repositories for $\mathcal{R}_{\mathrm{struct}}, \mathcal{R}_{\mathrm{sem}}, \mathcal{R}_{\mathrm{log}}$
  \For{$s \in S$} \Comment{Channel 1: deterministic extraction}
    \State extract structural and explicit semantic facts from $s$
    \State normalize them against $\mathcal{M}$
    \State register provenance and compilation targets
    \State admit resulting constraints into the appropriate repository
  \EndFor
  \For{$d \in D$} \Comment{Channel 2: LLM-assisted extraction}
    \State segment $d$ and retrieve relevant metamodel context
    \State $\mathcal{C}^{\mathrm{cand}} \gets \textsc{ExtractStructuredConstraints}(d,\mathcal{M},\mathcal{S}_c)$
    \For{$c \in \mathcal{C}^{\mathrm{cand}}$}
      \State compute $\alpha(c)$
      \If{$\alpha(c)=\bot$}
        \State quarantine $c$ for review; \textbf{continue}
      \EndIf
      \If{\textsc{Compatible}$(c,\mathcal{M})$ and \textsc{Compilable}$(c)$}
        \State admit $c$ with provenance, anchor, and backend metadata
      \Else
        \State quarantine $c$ with failure reason
      \EndIf
    \EndFor
  \EndFor
  \State derive dependency order $\Gamma$ from admitted constraints
  \State \Return $\mathcal{I}=(\mathcal{M},\mathcal{R},\mathcal{P},\Gamma)$
\EndFunction
\end{algorithmic}
\end{algorithm}


\begin{figure}[t]
  \centering
  \includegraphics[width=0.7\textwidth, trim=0.5cm 0.5cm 0.5cm 0.5cm, clip]{kg-model.pdf}
  \caption{Illustrative fragment of the ICM knowledge graph.}
  \Description{Illustrative knowledge-graph fragment used by the Integrated Constraint Model.
  Metamodel nodes represent classes, attributes, and abstract elements; constraint nodes attach structural or semantic obligations to those metamodel elements through typed links such as refersTo, hasAttribute, and TypeOf.}
  \label{fig:kg-model}
  \vspace{-0.5\baselineskip}
\end{figure}

\paragraph{Dynamic constraints.}
In addition to the persistent ICM, ATLAS allows per-request constraints that apply only to a single generation run. These are compiled into ephemeral L2 validators and executed together with the persistent validators, but they are not promoted into the long-lived repository unless they pass the same grounding and compatibility checks as ordinary Channel~2 constraints. This distinction prevents one-off user instructions from silently changing the persistent domain metamodel and constraint repository.

\paragraph{Compilation to executable enforcement.}
Admitted constraints are compiled into different executables depending on their family:
\begin{align}
\kappa_{\mathrm{struct}} &: \mathcal{R}_{\mathrm{struct}} \rightarrow \mathcal{A}_{\mathrm{prefix}}, \\
\kappa_{\mathrm{sem}} &: \mathcal{R}_{\mathrm{sem}} \rightarrow \mathcal{V}_{\mathrm{sem}}, \\
\kappa_{\mathrm{logic}} &: \mathcal{R}_{\mathrm{log}} \rightarrow \mathcal{V}_{\mathrm{log}}.
\end{align}
In the AUTOSAR implementation reported in this paper, $\mathcal{A}_{\mathrm{prefix}}$ is realized by JSON-Schema / Regex / GBNF-derived automata, $\mathcal{V}_{\mathrm{sem}}$ by SHACL-like graph validators, and $\mathcal{V}_{\mathrm{log}}$ by SMT-backed checks. More generally, however, Layer~2 may be realized by existing domain validators, rule engines, or commercial toolchains when those provide stronger or more complete coverage than a paper-specific formalization.

\begin{figure}[!htbp]
  \centering
  \includegraphics[width=0.7\linewidth]{unified-constraints.pdf}
  \vspace{-1\baselineskip}
  \caption[Stratified compilation with LLM fallback.]
  {Stratified compilation with LLM fallback and repository policy.
  \textbf{Left:} Channel 1 programmatic mapping compiles metamodel and ICM constraints 
  to L1 validators (JSON Schema, DFA, GBNF) and L2 validators (SHACL, SMT), 
  preserving semantics at the chosen backend.
  \textbf{Right:} when programmatic mapping is incomplete, an LLM-based 
  pipeline~\cite{maSpecGenAutomatedGeneration2025} synthesizes missing 
  constraints with solver-backed validation; only verified outputs are 
  persisted in the ICM storage.}
  \Description{Two-panel diagram of stratified compilation with LLM fallback and repository policy.
  Left panel: Channel 1 performs programmatic mapping from metamodel and ICM constraints into L1, which enforces constraints at generation time using JSON Schema, FSM, GBNF, Regex, or DFA, and into L2, which enforces constraints post-generation via SHACL and SMT, while preserving intended semantics.
  Right panel: when the programmatic mapping is incomplete, a SpecGen-style pipeline invokes an LLM with solver-backed validation to synthesize missing SHACL/SMT constraints. Only constraints that pass validation are written into the Formal Constraints Database (ICM storage).}
  \label{fig:unified-constraints}
\end{figure}

\begin{proposition}[Scoped Compilation Contract]
\label{prop:scoped-compilation-contract}
Fix a bounded unfolding depth $d$ for recursive structural fragments. Let $\kappa_{\mathrm{struct}}$, $\kappa_{\mathrm{sem}}$, and $\kappa_{\mathrm{logic}}$ compile admitted constraints into the corresponding L1 and L2 executables. Then ATLAS relies on the following scoped contract:
\begin{enumerate}[label=(\roman*),leftmargin=*,itemsep=0pt]
  \item L1 masking is applied only to structural constraints that have been compiled into prefix-decidable automata under the chosen bound $d$.
  \item Any draft accepted by the compiled L1 mechanism satisfies the compiled structural obligations represented by that executable controller.
  \item Any completed artifact accepted by all compiled L2 validators satisfies the admitted semantic and logical obligations represented by those validators.
\end{enumerate}
No claim is made for constraints that were not admitted into the ICM, that exceed the chosen unfolding bound, or that are only partially approximated by the backend encoding.
\end{proposition}

\paragraph{Rationale.}
Item (i) states the admission policy: only prefix-checkable structural fragments are sent to L1. Item (ii) therefore holds with respect to the compiled controller and the selected unfolding depth, while item (iii) holds with respect to the executable predicates realized by the validator backends.

\subsection{Unified Automaton Execution}
\label{sec:generation-strategy}

\paragraph{Rationale for the L1/L2 Split.}
Some obligations can be enforced while the artifact is still being written; others cannot. Local structural conditions such as nesting, token class, field presence, and ordering are prefix-checkable and can therefore be compiled into a decoder-side controller. Cross-reference integrity, graph acyclicity, type compatibility across distant elements, or numeric timing constraints require the completed artifact and belong in L2. This separation keeps decode-time control focused on prefix-checkable structure while reserving global semantic and logical conditions for post-generation validation.

\begin{figure}[!htbp]
  \centering
  \includegraphics[width=\linewidth, trim=0.5cm 0.5cm 0.5cm 0.5cm, clip]{guided_decode.pdf}
  \caption[Unified automaton execution for L1 enforcement during decoding.]
  {
  Unified automaton execution for L1 enforcement during LLM decoding.
  \textbf{Input constraints:} L1 structural constraints
  ($\mathcal{R}_{\mathrm{struct}}$ from the ICM; Section~\ref{sec:icm})
  are compiled into executable automata.
  Finite-state fragments (JSON Schema / Regex / FSM) are determinized to prefix-closed DFAs.
  GBNF/CFG fragments are executed via PDA/LR; for unified masking and logging, bounded-depth unfolding (depth $d$) produces an equivalent DFA $\mathcal{A}_d$.
  \textbf{Runtime decoding loop:} At decoding step $t$, the executor computes
  \(
  M_t = \{y \in \Sigma \mid \delta(s_t,y) \neq \bot\}
  \),
  masks invalid tokens in the LLM logits, samples a valid token $y_t$, and transitions to $s_{t+1} = \delta(s_t,y_t)$.
  \textbf{Audit trail:} Each step records
  $\tau_t = \langle s_t, |M_t|, y_t, s_{t+1}, \Delta t \rangle$
  into the audit trail.
  Post-generation, L2 validators check the completed artifact, yielding an evidence bundle $\Pi = \langle \rho_{\text{struct}}, \rho_{\text{sem}}, \rho_{\text{logic}}, \tau \rangle$ (Section~\ref{sec:evidence-bundle}).}
  \Description{Diagram of unified automaton execution for Layer 1 (L1) enforcement during large language model decoding.
  First, structural constraints from the Integrated Constraint Model (ICM) are compiled into executable automata: JSON Schema, Regex, and FSM fragments become deterministic finite automata, while GBNF and CFG fragments are executed via PDA/LR and optionally unfolded to bounded-depth DFAs for unified masking and logging.
  During decoding at each step t, the executor computes the set of valid next tokens, masks invalid tokens in the model logits, samples a valid token, and advances the automaton state.
  Each decoding step logs an audit tuple including states, mask size, chosen token, and timing.
  After generation finishes, Layer 2 (L2) validators for semantic constraints and logical constraints check the final artifact and combine the audit trail into an ordered evidence bundle.}
  \label{fig:guided_decode}
\end{figure}

Let $s_t$ denote the automaton configuration at decoding step $t$. The allowed token set is
\[
M_t
=
\{y \in \Sigma \mid \delta(s_t,y)\neq\bot \text{ and } \textsc{ReachAccept}(\delta(s_t,y))\},
\]
where $\textsc{ReachAccept}$ checks whether the successor configuration can still reach an accepting state under the compiled structural controller.

In practice, ATLAS uses DFA execution for finite-state fragments and bounded unfolding for recursive grammar fragments. The latter is a pragmatic compromise: it gives the decoder a uniform masking interface while making the scope of structural guarantees explicit through the bound $d$.

\paragraph{Prefix safety and structural closure.}
Under the L1 masking policy, a token is admitted only when the successor controller state can still reach acceptance in the compiled structural automaton. This preserves at least one admissible continuation for every emitted prefix. When decoding halts before acceptance, the controller can compute an accepting continuation on the compiled graph and use it as a structural closure mechanism for an incomplete draft. In our setting, this closure step is a controller-level completion procedure for restoring L1 admissibility; semantic and logic checks remain the responsibility of Layer~2 validation.

\paragraph{Runtime considerations.}
L1 enforcement changes the model's token distribution and can over-favor short valid completions. To counter this, ATLAS uses minimum-coverage guards for required regions, bounded unfolding calibration, and a two-stage workflow in which L1 focuses on structure while L2 and AGR address richer semantic issues. These are engineering controls rather than formal guarantees, but they are important for preserving content richness instead of generating the shortest acceptable artifact.

\subsection{Constraint-Guided, Validation-Backed Generation}
\label{sec:evidence-bundle}

\paragraph{Goal and scope.}
CVG produces an artifact together with an evidence bundle that can be re-checked independently of the generator. The resulting output is intended to support downstream review, validation, and repair.


\begin{figure}[!htbp]
  \centering
  \includegraphics[width=0.9\linewidth, trim=0.5cm 0.5cm 0.5cm 0.5cm, clip]{validation-evidence.pdf}
  \caption[Validation and evidence composition in ATLAS.]{ATLAS's validation and evidence composition. The generated artifact and the decode-time audit trail are checked by structural, semantic, and logic validators derived from the ICM. Each stage emits a validation record. These records are bundled into a single evidence package that is used both to accept or reject the artifact and to drive targeted repair when checks fail.}

  \Description{Workflow for validation and evidence composition in ATLAS.
  Inputs include the generated artifact, the runtime audit trail from Layer 1 guided decoding, and validators compiled from the Integrated Constraint Model (ICM): Layer 1 structural validators such as DFA, JSON Schema, and GBNF, and Layer 2 semantic and logical validators such as SHACL and SMT, including constraints produced with LLM assistance.
  Validation executes in structural, then semantic, then logic order. Each stage produces a validation record that is appended to the evidence bundle.
  These records are combined into an ordered bundle together with the audit trail. A bundle checker evaluates this package; if it passes, the artifact is exposed together with its validation bundle, and if it fails, an Audit-Guided Repair process uses the evidence to generate targeted fixes.}
  \label{fig:validation-evidence}
\end{figure}

\paragraph{Validation-backed evidence bundle.}
For each generated artifact $a$, ATLAS records an ordered evidence bundle
\[
\Pi=\langle \rho_{\text{struct}},\rho_{\text{sem}},\rho_{\text{logic}},\tau \rangle,
\]
where the entries summarize the outcomes of L1 structural checking, L2 semantic validation, L2 logic validation, and the decode-time audit trail, respectively.

The evidence bundle contains four parts:
\begin{itemize}[leftmargin=*,itemsep=0.4ex]
  \item $\rho_{\text{struct}}$: the structural acceptance record, including the identifier of the compiled controller and the outcome of L1 checking;
  \item $\rho_{\text{sem}}$: semantic validation reports, such as SHACL-style conformance results and localized violation paths;
  \item $\rho_{\text{logic}}$: logic validation results, such as solver satisfiability outcomes or conflict sets;
  \item $\tau$: the decode-time audit trail produced during L1-guided generation.
\end{itemize}

At decoding step $t$, ATLAS records
\[
\tau_t=\langle s_t,\ |M_t|,\ y_t,\ s_{t+1},\ \Delta t \rangle,
\]
where $s_t$ and $s_{t+1}$ are pre- and post-token controller configurations, $|M_t|$ summarizes mask restrictiveness, $y_t$ is the emitted token, and $\Delta t$ is the per-step latency. The full audit trail is the ordered sequence $\tau=\langle \tau_1,\ldots,\tau_T\rangle$.

We treat $\Pi$ as an ordered bundle rather than as a separate algebra: validator outcomes determine acceptance, while $\tau$ provides process metadata for replay, diagnosis, and repair prioritization. Operationally, an artifact is accepted only when the structural, semantic, and logic validators all pass. Audit metadata does not by itself turn a failed artifact into an accepted one; it explains \emph{how} generation proceeded and supports diagnosis when a validator fails.

\begin{table}[!htbp]
\centering
\small
\caption{Constraint hierarchy and evidence mapping.
Each constraint family is enforced at its natural layer and emits a corresponding audit/validation artifact.}
\label{tab:hier-matrix}
\resizebox{\linewidth}{!}{%
\begin{tabular}{@{}p{0.16\columnwidth}p{0.26\columnwidth}p{0.29\columnwidth}p{0.25\columnwidth}@{}}
\toprule
\textbf{Type} &
\textbf{Layer-1 (Generation Time)} &
\textbf{Layer-2 (Post Generation)} &
\textbf{Evidence Artifact} \\
\midrule
Structural
& JSON Schema / GBNF $\rightarrow$ automaton masking
& structural re-check and closure summary
& $\rho_{\text{struct}}$, controller id, audit trail \\
Semantic
& local enums / local field guards when prefix-checkable
& graph-level conformance and reference checks
& $\rho_{\text{sem}}$ with violation paths and anchors \\
Logic
& local numeric guards when safely compilable
& solver-backed numeric / temporal checks
& $\rho_{\text{logic}}$ with models or conflict sets \\
\bottomrule
\end{tabular}
}
\end{table}

Table~\ref{tab:hier-matrix} makes one design choice explicit: some rules can appear in both layers in different forms. A simple enumeration may be enforced during decoding, while a richer compatibility obligation involving multiple elements is checked afterward. This redundancy allows ATLAS to block simple structural faults during decoding while reserving richer domain reasoning for post-generation validation.

\subsection{Audit-Guided Intelligent Repair}
\label{sec:validation-repair}

\paragraph{Role in the pipeline.}
By the time an artifact reaches AGR, structural validity has already been enforced by L1 and summarized in $\rho_{\text{struct}}$. Remaining failures are typically semantic or logical: broken references, incompatible types, acyclicity violations, timing conflicts, or domain-specific numeric inconsistencies. AGR therefore focuses on evidence-driven correction rather than free-form regeneration.

AGR consumes the ordered evidence bundle
\[
\Pi
=
\langle \rho_{\text{struct}}, \rho_{\text{sem}}, \rho_{\text{logic}}, \tau \rangle
\]
and turns validator feedback into localized repair tasks. The central idea is that validators should not only reject artifacts; they should emit failure objects that identify \emph{where} the problem is, \emph{which obligation} was violated, and \emph{which metamodel or ICM item} that obligation came from.

\paragraph{Violation localization.}
For semantic failures, AGR uses validator outputs such as entity identifiers, paths, and expected relation types to locate the failing region in the artifact and recover the associated metamodel anchor. For logic failures, it uses solver-level diagnostics such as violated clauses or unsat cores to identify the conflicting fields and obligations. These localized diagnostics are crucial because they let ATLAS repair specific regions while preserving the already accepted parts of the artifact.

\paragraph{Two repair routes.}
ATLAS supports two repair routes:
\begin{enumerate}[label=(\roman*),leftmargin=*,itemsep=0.5ex]
  \item Automated validator-guided patching.
        An LLM receives the localized diagnostics, the relevant metamodel context, and the violated constraint text, and is asked to produce a minimal patch for the affected region.
  \item Human-in-the-loop typed editing.
        When the violation is high-stakes, ambiguous, or repeatedly unstable under automated patching, the artifact is projected into a metamodel-derived editor that lets a domain expert correct the typed model directly while seeing the same validator diagnostics.
\end{enumerate}

This division of repair routes reflects the intended automation boundary of ATLAS: automated patching addresses structurally regular or locally diagnosable failures, while human intervention is reserved for ambiguous or high-stakes cases that require domain judgment.

\paragraph{Constraint promotion.}
When a repair reveals a reusable obligation that was previously absent or too weakly enforced, ATLAS can promote that rule back into future generation contexts: for example by augmenting retrieval snippets, tightening a local L1 guard, or persisting a new L2 validator. Promotion is gated by the same provenance and compatibility checks used during initial ICM construction, so repair does not become an uncontrolled source of new global rules.

\paragraph{Dependency-aware scheduling.}
AGR orders fixes using the dependency relation $\Gamma$ from the ICM. In practice, this means addressing prerequisite faults first: unresolved references before downstream type checks, and type or cardinality mismatches before numeric tuning. This ordering makes repair more stable because later validators are not asked to reason over obviously malformed intermediate states. The helper procedures in Algorithm~\ref{alg:repair} are likewise workflow-level interfaces for violation extraction, local revalidation, and patch application rather than a claim of a closed repair calculus.

\begin{algorithm}[t]
\caption{Layer-2-Guided Repair with Human/LLM Routes and Constraint Promotion}
\label{alg:repair}
\begin{algorithmic}[1]
\small
\Require Artifact $a$, evidence bundle $\Pi$, ICM $\mathcal{I}$
\Ensure Repaired artifact $a^*$ or \texttt{MANUAL\_REVIEW}

\State $\mathcal{V}_{\text{sem}} \gets \textsc{ExtractSemanticViolations}(\rho_{\text{sem}})$
\State $\mathcal{V}_{\text{logic}} \gets \textsc{ExtractLogicViolations}(\rho_{\text{logic}})$
\State $\mathcal{V} \gets \mathcal{V}_{\text{sem}} \cup \mathcal{V}_{\text{logic}}$
\State order $\mathcal{V}$ using the dependency relation $\Gamma$ in $\mathcal{I}$

\For{$v \in \mathcal{V}$}
    \State recover provenance, anchor, and relevant metamodel context for $v$
    \If{\textsc{AutoRepairable}$(v)$}
        \State build a constrained repair prompt from diagnostics and context
        \State $a \gets \textsc{PatchWithLLM}(a,v)$
    \Else
        \State $a \gets \textsc{ApplyTypedEditorFix}(a,v)$
    \EndIf
    \State locally revalidate the affected region
    \If{revalidation fails repeatedly}
        \State \Return \texttt{MANUAL\_REVIEW}
    \EndIf
    \If{$v$ yields a reusable new rule}
        \State \textsc{PromoteConstraint}(v,$\mathcal{I}$)
    \EndIf
\EndFor
\State \Return $a$ as $a^*$
\end{algorithmic}
\end{algorithm}

In this paper, AGR provides the interface that connects validation evidence, localized repair, and human intervention, and it gives the evaluation a concrete basis for discussing remediation cost and remaining manual effort.


\section{Evaluation}
\label{sec:evaluation}
\subsection{Overview and Research Questions}
\label{sec:evaluation-overview}
%======================================================================

We evaluate ATLAS in AUTOSAR, a mature schema-rich MDE setting with authoritative schemas, cross-file references, and validator-backed semantics. This setting allows us to assess both what ATLAS can enforce during generation (structural validity) and what remains subject to post-generation diagnosis (global semantic and logical consistency).

We structure the evaluation around two questions:

\noindent \textbf{RQ1 (Single-File Structural Reliability):} Can Layer-1 constraint enforcement reliably produce structurally valid single-file AUTOSAR artifacts, and what latency and token overhead does this introduce?

\noindent \textbf{RQ2 (System-Level Scaling and Residual Failure Modes):} When moving to multi-file AUTOSAR systems, which properties remain robust under layered generation, which semantic and logical failures remain, and how visible are those failures to validator-guided repair and result analysis?

Evaluation combines automated validation (XSD, SHACL, SMT, reference-resolution, cost) with a diagnostic analysis of multi-file AUTOSAR outputs. This analysis characterizes residual defects, repair burden, and the boundary between machine-resolvable and analysis-identified issues. All experiments use fixed seeds or fixed API settings, documented model versions, and archived artifacts for replay and inspection.


\subsection{AUTOSAR Domain Preparation}
\label{sec:autosar-domain-setup}

ATLAS evaluates the authoritative-schema AUTOSAR setting. For the experiments, we construct three coupled assets that are reused throughout generation, validation, and repair: an integrated AUTOSAR metamodel, an AUTOSAR-specific ICM, and a queryable structural substrate derived from them. This preparation instantiates the framework with reproducible, machine-readable AUTOSAR resources.

First, we derive the integrated AUTOSAR metamodel deterministically from official AUTOSAR XSD and XMI artifacts, corresponding to the S1 pathway in Section~\ref{sec:methodology}. The transformation normalizes core AUTOSAR types---including software components, ports, interfaces, runnables, signals, and timing attributes---together with their containment, cardinality, and reference relations. This gives the pipeline a canonical vocabulary for downstream retrieval, constraint compilation, and validation.

Second, we construct the AUTOSAR ICM by linking that metamodel to constraints extracted from AUTOSAR PDF specifications and related normative documents through the dual-channel process described in Section~\ref{sec:icm}. The resulting constraint set captures a substantial machine-compilable subset of mandatory element presence, value-range restrictions, reference-integrity obligations, and timing or behavioral conditions in a validation-ready form. Applying this process to the AUTOSAR Software Component Template specification yielded \textbf{1,161 normative constraints}, of which 1,045 were anchored to verified metamodel entities (90.0\% anchoring rate). In the present AUTOSAR instantiation, these admitted constraints are compiled into \textbf{794 SHACL node shapes} for graph-structural checks and \textbf{1,005 SMT assertions} for logic and value checks; purely informational hints are excluded from the automated pipeline. These SHACL/SMT artifacts provide a concrete Layer~2 instantiation over a publicly describable subset of AUTOSAR obligations and make post-generation validation executable in the academic evaluation. The broader industrial workflow remains compatible with richer external AUTOSAR validators.

Third, we operationalize the metamodel and ICM as a queryable structural substrate used in two ways during evaluation: to assemble ICM-complete retrieval prompts for generation, and to provide the schema fragments and typed obligations needed by Layer~1 and Layer~2 validation. In practice, this substrate exposes the same AUTOSAR entities, relations, and attached constraints in a form that can be retrieved component-wise without manual per-instance rule rewriting. Overall, the AUTOSAR preparation stage yields a reproducible three-part asset base---metamodel, ICM, and operational retrieval/validation substrate---that underpins all experiments reported below.

\subsection{RQ1: Single-File AUTOSAR Component Generation}
\label{sec:rq1-autosar}
%======================================================================
RQ1 evaluates the complete ATLAS pipeline on single-component AUTOSAR generation, examining whether metamodel- and ICM-guided retrieval, layered constraint enforcement, and audit instrumentation deliver structural correctness with tractable decoding latency and token overhead.


\subsubsection{Experimental Setup}
\label{sec:rq1-setup}

\paragraph{Research Question.} RQ1 studies structural reliability in single-file AUTOSAR generation under two complementary forms of control: (i) ICM-complete retrieval prompting, which exposes deterministically stored metamodel structure and schema-compatible cues to the model at prompt time, and (ii) Layer-1 constrained decoding, which enforces admissible prefixes during generation. The question is therefore not only whether structurally valid artifacts can be obtained, but how prompt-side structural steering and decode-time admissibility control interact in terms of validity, auditability, latency, and token cost.

\paragraph{Baseline Comparisons.} We evaluate four \emph{controlled local} pipeline configurations that represent progressively structured approaches from the recent LLM literature, each implemented using the same base model to isolate the distinct effects of prompt-side structural steering and decode-time constraint enforcement:

\begin{enumerate}[leftmargin=*,nosep]
\item \textbf{vLLM} (pure prompting baseline): Direct generation without retrieval augmentation or Layer-1 enforcement, mirroring common practice in LLM-based information extraction systems where pretrained models are directly prompted to emit structured records~\cite{Dagdelen2024, Lu2022}.

\item \textbf{vLLM+ICM-RAG}: Augments prompts with an ICM-complete retrieval context assembled from deterministically stored AUTOSAR metamodel facts, typing relations, containment structure, and schema-compatible field cues. Its role is not decode-time enforcement, but structural steering: it conditions the model on an authoritative and structurally complete representation before generation~\cite{el-gnainyAIEnhancedAUTOSARConfiguration2024a,Jiao2024}.

\item \textbf{vLLM+ICM-RAG+JSON Schema}: Builds on the same ICM-complete prompting context, but adds Layer-1 constrained decoding through JSON-structured intermediate generation followed by deterministic projection to ARXML. This path turns structurally guided prompting into schema-level structural guarantees with replayable audit traces~\cite{Lu2025b}.

\item \textbf{vLLM+ICM-RAG+GBNF}: Also builds on the same ICM-complete prompting context, but realizes Layer-1 enforcement through grammar-constrained decoding with GBNF-derived masking over admissible prefixes, following the constrained decoding paradigm~\cite{scholakPICARDParsingIncrementally2021,Dong2024XGrammar,Park2025}.
\end{enumerate}


\paragraph{Prompt-side steering versus decode-time enforcement.}
ATLAS evaluates two complementary forms of structural control. The first is \emph{ICM-complete prompting}, which retrieves a structurally complete context from the ICM and exposes authoritative metamodel facts, typing relations, containment structure, and schema-compatible cues before decoding. The second is \emph{Layer-1 constrained decoding}, which enforces admissible prefixes during generation through executable structural contracts such as JSON Schema or GBNF.

This distinction matters because constrained decoding is known to introduce \emph{distribution distortion}: hard grammar masking can bias generation toward shorter but valid completions by eliminating longer yet still admissible continuations, thereby distorting the base model's learned distribution~\cite{parkGrammaralignedDecoding2025}. Grammar-Aligned Decoding (GAD) mitigates this effect through adaptive sampling with approximate expected futures, preserving grammaticality while better aligning with the base model's preferences~\cite{parkGrammaralignedDecoding2025}. However, such forward simulation introduces additional overhead that grows with grammar complexity and generation length, which can be undesirable for large industrial artifacts under deeply nested schemas. ATLAS therefore adopts a stratified design: ICM-complete prompting first increases the probability mass of schema-consistent continuations, while Layer-1 constrained decoding provides deterministic structural guarantees when strict structural correctness is required.

\paragraph{Model Choice and Scope.} RQ1 is designed as a \emph{mechanism-isolation study}: its purpose is to measure how much structural reliability comes from ATLAS itself---ICM-complete prompting, Layer-1 constrained decoding, and audit instrumentation---rather than from switching among model families. We therefore fix the backbone to DeepSeek-R1-Distill-Qwen-32B served through vLLM and treat it as a \emph{controlled constant}. This choice is methodological rather than doctrinal. An open-weight local model gives us a fully instrumented inference stack for tokenizer-level masking, deterministic replay, latency/token measurement, and stepwise audit capture, while remaining strong enough for differences among prompting, retrieval, and constrained decoding to remain visible. The point is not that this backbone is universally optimal, but that it lets RQ1 attribute observed differences to structural control mechanisms rather than to cross-model capability variance. Porting ATLAS to another backbone requires only three localized adjustments: a tokenizer-compatible Layer-1 masking interface, structured-output serving support, and retuning of sampling/context-budget parameters. The metamodel-integration stage, ICM, Layer-2 validators, and audit/repair logic are otherwise unchanged.

The four controlled local pipelines use the same backend model (DeepSeek-R1-Distill-Qwen-32B served through vLLM) with fixed random seeds (42, 1001, 20250701), temperature 0.7, and top-p sampling at 0.9 to ensure reproducibility. We evaluate 60 representative AUTOSAR components under three prompt regimes that vary the amount of natural-language specification context provided: \textsc{Min} (minimal requirements), \textsc{Std} (standard documentation), and \textsc{Full} (comprehensive specifications including timing and behavioral constraints). In addition, Table~\ref{tab:rq1-correctness} reports one \emph{external reference baseline} using GPT-5 in an unconstrained API setting. This API baseline serves a different purpose from the controlled local pipelines: it measures how far a strong commercial model can go without ATLAS-style structural steering or Layer-1 enforcement. It is therefore included only as a supplementary frontier reference for syntax feasibility and is not part of the same-backbone latency or audit comparison.

\paragraph{Evaluation Metrics.} For this structural feasibility study, we focus on three dimensions:

\begin{itemize}[leftmargin=*,itemsep=0.5ex]
\item \textbf{Structural correctness}: XSD validation pass rate indicating syntactic compliance with AUTOSAR schemas.
\item \textbf{Audit completeness}: Audit trail completeness measuring the percentage of generation steps for which per-step audit events $\tau_t$ are recorded; we report structural acceptance record availability separately at the artifact level via $\rho_{\text{struct}}$.
\item \textbf{Computational cost}: End-to-end latency in seconds and token counts for input and output.
\end{itemize}

\textit{Note: Semantic consistency (SHACL) and logic verification (SMT) metrics are reserved for the system-level evaluation in RQ2, as single-file generation lacks the necessary cross-reference context.}


\subsubsection{Results}
\label{sec:rq1-results}

\paragraph{Computational Cost.} Table~\ref{tab:rq1-latency} presents end-to-end latency and token consumption across pipelines and prompt regimes. The pure vLLM baseline exhibits moderate latency and minimal input tokens due to lack of retrieval augmentation, but generates longer outputs containing structural errors. Adding ICM-complete retrieval increases input context substantially (from 154--227 tokens to 992--6,649 tokens depending on regime) while improving output quality by supplying structurally complete context. Constrained decoding pipelines (GBNF and JSON Schema) reduce output tokens by enforcing compact, schema-compliant generation, with JSON Schema achieving the lowest latency due to deterministic intermediate representation.

\begin{table}[!htbp]
\centering
\small
\caption{Computational cost for RQ1: end-to-end latency and token counts by pipeline and prompt regime. All measurements use DeepSeek-R1-Distill-Qwen-32B on the same hardware.}
\label{tab:rq1-latency}
\begin{tabular}{@{}lccc@{}}
\toprule
\textbf{Pipeline Configuration} & \textbf{Latency (seconds)} & \textbf{Input Tokens} & \textbf{Output Tokens} \\
\midrule
\texttt{vLLM} (\textsc{Min})      & 21.87 & 154  & 812  \\
\texttt{vLLM} (\textsc{Std})      & 29.41 & 188  & 1,279 \\
\texttt{vLLM} (\textsc{Full})     & 46.31 & 227  & 2,013 \\
\midrule
\texttt{vLLM+ICM-RAG} (\textsc{Min})  & 17.51 & 992  & 743  \\
\texttt{vLLM+ICM-RAG} (\textsc{Std})  & 33.19 & 2,534 & 1,397 \\
\texttt{vLLM+ICM-RAG} (\textsc{Full}) & 85.60 & 4,493 & 3,610 \\
\midrule
\texttt{vLLM+ICM-RAG+GBNF} (\textsc{Min})  & 8.01   & 828  & 155  \\
\texttt{vLLM+ICM-RAG+GBNF} (\textsc{Std})  & 59.94  & 2,896 & 1,446 \\
\texttt{vLLM+ICM-RAG+GBNF} (\textsc{Full}) & 123.74 & 6,649 & 3,118 \\
\midrule
\texttt{vLLM+ICM-RAG+JSON Schema} (\textsc{Min})  & 3.83  & 848  & 119  \\
\texttt{vLLM+ICM-RAG+JSON Schema} (\textsc{Std})  & 13.38 & 2,605 & 523  \\
\texttt{vLLM+ICM-RAG+JSON Schema} (\textsc{Full}) & 31.56 & 2,915 & 1,291 \\
\bottomrule
\end{tabular}
\end{table}

\paragraph{Audit Completeness.} For the vLLM+ICM-RAG+JSON Schema pipeline, the structured-output path records complete per-step audit coverage with Layer-1 candidate sets and also emits a structural acceptance record $\rho_{\text{struct}}$ for all 60 test cases across all three prompt regimes. Table~\ref{tab:rq1-audit} summarizes these audit statistics. The allowed token set size averages 156 tokens per generation step with minimum of 1 (fully constrained choices) and maximum of 650 (unconstrained text fields), demonstrating effective constraint propagation through the controller. Audit trace density approximates three recorded events per accepted token, providing fine-grained provenance for verification. The vLLM+ICM-RAG+GBNF pipeline was not instrumented with the same audit recorder in our current implementation, representing an engineering limitation rather than fundamental incompatibility with GBNF-based constraint enforcement.

\begin{table}[!htbp]
\centering
\caption{Layer-1 auditing statistics for vLLM+ICM-RAG+JSON Schema pipeline, showing audit coverage and constraint effectiveness across prompt regimes.}
\label{tab:rq1-audit}

\resizebox{\linewidth}{!}{%
\begin{tabular}{@{}lcccc@{}}
\toprule
\textbf{Prompt Regime} &
\textbf{Allowed Token Set (median / min)} &
\textbf{Generation Steps} &
\textbf{Audit Coverage} &
\textbf{Structural Acceptance Record Availability} \\
\midrule
\textsc{Min}  & 156 / 1 (max $\leq$ 650) & 119   & 100\% & 100\% \\
\textsc{Std}  & 156 / 1 (max $\leq$ 650) & 524   & 100\% & 100\% \\
\textsc{Full} & 156 / 1 (max $\leq$ 650) & 1{,}292 & 100\% & 100\% \\
\bottomrule
\end{tabular}%
}
\end{table}


\paragraph{Structural Correctness (Syntax Feasibility).}
Since single-component generation in isolation lacks the system-level context required to satisfy global semantic constraints, such as references to external interfaces, RQ1 focuses strictly on structural validity, which is the prerequisite for any downstream engineering utility.
Table~\ref{tab:rq1-correctness} presents the XSD validation pass rates.
The results reveal three distinct regimes. The pure vLLM baseline failed completely (0\% pass rate), indicating that unconstrained prompting alone could not reliably satisfy the strict nesting and typing discipline of AUTOSAR ARXML. The frontier LLM-API baseline (GPT-5), despite substantially stronger general-purpose capability, achieved only a 50\% pass rate when used without ATLAS-style structural control.

By contrast, \texttt{vLLM+ICM-RAG} achieved 100\% XSD validity. This should be interpreted as evidence for the effectiveness of \emph{ICM-complete prompting} rather than generic retrieval alone: the retrieved context deterministically exposes metamodel structure and schema-compatible cues, allowing the model to operate under a structurally complete prompt. The Layer-1-constrained pipelines (\texttt{vLLM+ICM-RAG+JSON Schema} and \texttt{vLLM+ICM-RAG+GBNF}) matched this validity while strengthening it into decode-time admissibility control. In other words, ICM-complete prompting makes valid continuations more probable, whereas constrained decoding makes invalid continuations unreachable with respect to the compiled structural contract.

\begin{table}[!htbp]
  \centering
  \small
  \caption{Structural correctness for RQ1: XSD pass rate per pipeline.}
  \label{tab:rq1-correctness}
  \begin{tabular}{@{}lc@{}}
    \toprule
    \textbf{Method}                 & \textbf{XSD Pass Rate (\%)} \\
    \midrule
    \texttt{vLLM}                   & 0   \\
    \texttt{vLLM+ICM-RAG}               & 100 \\
    \texttt{vLLM+ICM-RAG+GBNF}          & 100 \\
    \texttt{vLLM+ICM-RAG+JSON Schema}   & 100 \\
    \texttt{LLM-API}                & 50  \\
    \bottomrule
  \end{tabular}
\end{table}


\subsubsection{Analysis and Discussion}
\label{sec:rq1-analysis}

\paragraph{Validation of Layered Constraint Architecture.}
RQ1 shows that structural reliability in ATLAS is achieved through two complementary mechanisms rather than a single intervention. The first is \emph{ICM-complete prompting}, which supplies a structurally complete and authoritative context before decoding. The second is \emph{Layer-1 constrained decoding}, which enforces the same structural contract during generation. Their roles are different. Prompt-side grounding improves distributional alignment with schema-consistent continuations, whereas decode-time masking provides explicit structural guarantees and, in the JSON-Schema path, complete replayable audit traces.

This distinction helps explain why \texttt{vLLM+ICM-RAG} already attains perfect XSD validity in this benchmark. The result does not imply that ordinary vector-only RAG would do the same; rather, it shows that when retrieval is anchored in the ICM, the prompt itself becomes a strong structural prior. Adding Layer-1 constraints then shifts the guarantee from empirical success under a structurally complete prompt to explicit structural enforcement at decode time.

The \texttt{LLM-API} result further clarifies the point. Greater model capability alone did not substitute for structural control: without ATLAS-style grounding and enforcement, even a frontier commercial model remained unreliable under AUTOSAR's schema discipline. By combining probability shaping at the prompt level with structural enforcement during decoding, ATLAS turns the LLM into a structurally disciplined artifact generator whose accepted outputs remain ingestible by downstream engineering tools.

\paragraph{Audit-Backed Generation and Trace Visibility.}
The 100\% audit coverage and availability of structural acceptance records for the instrumented constrained pipeline show that audit-backed constrained generation is practically realizable without prohibitive overhead.
The audit trail density provides sufficient granularity to trace structural decisions back to the enforcing automaton state.
This recorded trace enables fine-grained replay and inspection of decode-time enforcement decisions, supporting reproducibility and auditability in assurance-oriented engineering workflows—a capability absent in standard API-based generation.

\paragraph{Latency and Token Cost.}
The latency and token cost results (Table~\ref{tab:rq1-latency}) reveal the trade-offs of constrained generation. Constraint-aware retrieval substantially increases input tokens (up to 6,649 tokens in the \textsc{Full} regime) to provide the necessary context, while reducing output tokens through more compact and compliant generation. The JSON Schema pipeline achieves the lowest latency by using deterministic intermediate generation followed by deterministic JSON-to-ARXML serialization; the reported latency includes this serialization step. GBNF exhibits higher latency due to the overhead of per-token automaton state transitions. These metrics suggest that JSON Schema-based enforcement offers the strongest balance between performance and structural reliability for this domain.

\paragraph{Synthesis: Structural Validity as a Foundation.}
RQ1 shows that ATLAS's prompt-side grounding and Layer-1 enforcement together address the ``syntax problem'' of LLM-based generation in this benchmark.
While single-file generation inherently precludes verification of global cross-file semantics (the subject of RQ2), the observed 100\% structural validity remains important because it ensures that downstream semantic validation operates on well-formed, parseable artifacts.
Without this structural control, higher-level semantic reasoning and repair would first have to absorb basic parsing failures, as evidenced by the unconstrained baselines.
RQ1 therefore establishes the structural foundation on which the multi-file semantic assembly studied in RQ2 can be evaluated.
% ------------------------------------------------------------------------------------------------------------------------
%======================================================================
\subsection{RQ2: Multi-File System Generation at Scale}
\label{sec:rq2-multifile-autosar}
%======================================================================

The second research question extends the evaluation from the single-component setting of RQ1 to system-level, multi-file AUTOSAR configurations. Here we stress-test ATLAS under realistic cross-file dependencies, asking whether the ATLAS architecture can maintain the strong Layer-1 structural reliability observed in RQ1 while surfacing the semantic and logical failure modes that only emerge at system scale.


\subsubsection{Experimental Setup}
\label{sec:rq2-setup}

\begin{figure}[!htbp]
\centering
\includegraphics[width=0.7\textwidth, trim=0.5cm 0.5cm 0.5cm 0.5cm, clip]{autosar-exp-arch.pdf}%
\vspace{-0.5\baselineskip}
\caption[Two-phase AUTOSAR generation, validation, and repair pipeline.]
{Two-phase AUTOSAR generation, validation, and repair pipeline. Phase~1 
synthesizes a high-level blueprint from NLP requirements and AUTOSAR standards, 
populating the metamodel and ICM substrate. Phase~2 performs constrained LLM decoding to 
generate component-specific ARXML with JSON Schema enforcement. Post-generation 
validation applies XSD, SHACL, and SMT checks, with Audit-Guided Repair 
generating targeted fixes for any violations to produce a validated AUTOSAR 
artifact bundle.}
\Description{System architecture for the AUTOSAR experiment shown as a
two-phase pipeline. Phase 1 (AUTOSAR Blueprint) consumes NLP requirements,
AUTOSAR standards/PNFs, and the AUTOSAR meta-model (XSD/XMI) to synthesize
a high-level component and type layout, populating the metamodel and ICM substrate and
the AUTOSAR knowledge graph. Phase 2 (Component-wise Generation) assembles
a dynamic JSON Schema for each component, performs constrained LLM decoding
to emit per-component JSON, and converts that JSON to ARXML. The workflow
retains phase-level provenance and generated artifacts for later inspection.
After generation, the produced ARXML is checked by XSD validation, SHACL semantic
validation, and SMT-based logical checks. A Failure Mode Logger and
Audit-Guided Repair module analyze any violations to generate targeted
fixes, yielding a validated AUTOSAR artifact bundle for downstream engineering
review and tool import.}
\label{fig:autosar-exp-arch}
\vspace{-0.5\baselineskip}
\end{figure}

\paragraph{Research Question and Motivation.} RQ2 tests whether ATLAS can scale to multi-file AUTOSAR systems while preserving Layer-1 structural correctness, and to what extent Layer-2 semantic and logical constraints can be satisfied in the presence of thousands of cross-file references and global invariants that do not arise in RQ1.


\paragraph{Dual-Phase Pipeline Architecture.} To address these challenges, we adopt a blueprint-guided dual-phase generation strategy that decomposes system-level generation into two stages, each with distinct objectives and constraint enforcement mechanisms:

\begin{enumerate}[leftmargin=*,nosep]
\item \textbf{Phase 1 (Blueprint Synthesis):} The system first generates a JSON-structured blueprint specifying the component inventory, dependency graph, file layout, and cross-component interface contracts. This blueprint serves as a global coordination artifact that establishes naming conventions, reference targets, and architectural topology before detailed component generation begins. The blueprint generation employs JSON Schema constraints to ensure structural validity and completeness of the system plan.

\item \textbf{Phase 2 (Component Assembly):} Guided by the blueprint, the system generates individual ARXML files for each component, with interfaces emitted as separate files referencing a shared basic-types definition. Each generation step uses the blueprint as context to resolve cross-component dependencies and maintain reference consistency. Component-level generation employs the same JSON Schema-constrained decoding used in Phase 1, ensuring structural compliance at the artifact level.
\end{enumerate}

This dual-phase decomposition addresses scalability through separation of concerns: Phase 1 handles global architectural decisions that require system-wide reasoning, while Phase 2 focuses on component-level detail generation constrained by the established blueprint. The blueprint acts as a contract between phases, reducing the context window requirements for Phase 2 generation while providing explicit targets for cross-file references.

\paragraph{Model Selection and Constraint Enforcement.} RQ2 serves a different evidential role from RQ1. Whereas RQ1 fixes a local open-weight backbone to isolate the contribution of ATLAS control mechanisms, RQ2 asks a boundary question: when raw model capability is increased in a deployment-oriented setting, which defects still remain at multi-file system scale, and which of them are still exposed by the layered ATLAS workflow? For that reason, RQ2 uses a frontier API-based model (GPT-5) in both phases. The goal is not a same-backbone comparison with RQ1, but a stress test of ATLAS under a stronger and more deployment-relevant model setting. This choice reflects three considerations:

\begin{itemize}[leftmargin=*,itemsep=0.5ex]
\item \textbf{System-scale reasoning demands}: Blueprint synthesis for multi-file AUTOSAR systems requires planning over component inventories, dependency graphs, interface contracts, and naming consistency before file-level realization begins. We therefore choose a stronger frontier model to test whether residual failures persist even when the underlying model has substantially greater general reasoning capacity.

\item \textbf{Deployment realism}: Many organizations consume advanced LLM capability through API access rather than a self-hosted inference stack, especially for occasional high-value synthesis tasks. RQ2 therefore complements the controlled local study of RQ1 with a realistic deployment setting instead of simply pushing the local 32B configuration to its upper bound.

\item \textbf{ATLAS continuity under API serving}: Although API-based generation precludes the token-level audit provenance demonstrated in RQ1, both phases still preserve Layer-1 structural control through the provider's JSON Schema structured-output interface. Phase 1 blueprint generation uses a JSON Schema specifying component inventory structure, dependency graph format, and file layout constraints. Phase 2 component generation uses JSON Schemas derived from AUTOSAR metamodel fragments, ensuring that individual ARXML files satisfy structural requirements before deterministic XML serialization. Thus, what changes in RQ2 is the serving environment and audit granularity, not the ATLAS principle that structural validity should be enforced during generation.
\end{itemize}

To balance creative architectural exploration with structural conformance, we adopt stage-specific temperature settings: Phase 1 blueprint synthesis uses temperature 0.7 to encourage diverse architectural solutions, while Phase 2 detailed generation uses temperature 0.3 for more deterministic output, and interface file emission uses temperature 0.2 for maximum consistency. 

\paragraph{Result Analysis Scope.}
The result analysis in RQ2 is diagnostic: it assesses whether generated systems preserve architectural intent, the severity of residual defects, and the classes of issues that still require correction after layered generation and validation. It is used to characterize residual defect burden and repair needs at system scale. The SHACL/SMT backends used in RQ2 provide reproducible Layer~2 instantiations for the AUTOSAR case study.

\paragraph{Dataset and Evaluation Metrics.} We curate 20 AUTOSAR systems across three complexity tiers designed to stress different aspects of cross-file generation. The dataset spans three distinct architectural configurations that progressively challenge the framework's ability to maintain cross-file consistency and semantic correctness.

The Simple tier comprises 7 cases with 27 components total, featuring chain topologies suitable for sensor-processor-actuator pipelines. Representative scenarios include data collection systems, control loops, and monitoring applications. Each system contains an average of 3 to 4 components, generating 52 ARXML files across the tier. The Middle tier consists of 7 cases with 45 components total, employing multi-branch tree topologies that represent sensor fusion architectures and multi-domain controllers integrating powertrain, body, and ADAS subsystems. Systems in this tier average 6 to 7 components each, producing 108 ARXML files. The Complex tier includes 6 cases with 54 components total, utilizing mesh topologies with bidirectional dependencies that model sophisticated scenarios such as autonomous driving stacks, by-wire chassis control, and vehicle-cloud coordination. These systems average 8 to 10 components each and generate 124 ARXML files.

We evaluate generated systems across six primary dimensions. File completeness measures whether all blueprint-declared files are produced without truncation, assessed post-generation for all 284 files across the 20 systems as a binary complete-or-incomplete determination. Structural correctness quantifies syntactic compliance with AUTOSAR XSD schemas through xmllint validation, reporting the XSD validation pass rate for all generated ARXML files. Cross-file reference integrity calculates the resolution rate as the proportion of resolvable references over total \texttt{<REF>} elements, measured by parsing all reference targets and verifying their existence in the generated artifact set. Toolchain usability evaluates whether systems successfully import into commercial AUTOSAR configuration tools, specifically Etas (ISOLAR-A), yielding a binary pass-or-fail outcome per system. Generation cost captures both token counts and wall-clock time, measured separately for Phase 1 blueprint synthesis and Phase 2 component assembly across all 20 systems.

Result analysis follows a structured protocol evaluating four dimensions: requirements alignment assesses whether generated systems satisfy functional specifications and maintain traceability to original requirements; architectural quality evaluates dependency consistency and type-level correctness; engineering quality measures naming conventions, structural modularity, and estimated repair effort; toolchain integration validates whether systems import successfully into commercial AUTOSAR configuration environments. For readability, Table~\ref{tab:rq2-analysis} uses the following rubric codes: A1 = functional coverage, A2 = interface completeness, A3 = requirements traceability; B1 = dependency/role consistency, B2 = local type correctness; C1 = naming/documentation quality, C2 = structural modularity, C3 = estimated repair effort; D1 = toolchain import success.

\subsubsection{Results}
\label{sec:rq2-results}

\paragraph{Structural Success vs. Semantic Failure.}
Table~\ref{tab:rq2-structural} (updated with semantic metrics) illustrates a distinct degradation in model performance as constraints shift from syntactic to semantic and from local to global.
At Layer~1, all 20 systems across all complexity tiers achieved 100\% XSD pass rates, confirming that blueprint-guided constrained decoding effectively solves the syntax-generation problem.
At Layer~2, local semantics checked via SHACL degrade sharply. Among the instantiated checks for the generated systems---57, 113, and 63 effective checks across the Simple, Middle, and Complex tiers, drawn from a validator inventory of 794 node shapes---the pass rates fall to 2/57 (about 3.5\%), 16/113 (about 14.2\%), and 6/63 (about 9.5\%), respectively. These failures largely stem from missing mandatory attributes or incorrect local configurations that remain structurally valid but semantically vacuous.
When SMT constraints are added to verify cross-file consistency, such as end-to-end signal flow and cross-component consistency obligations, the system-level pass rate drops to 0\% across all tiers.
While the model often generated references that resolved to \emph{existing} IDs, as shown by the high reference-resolution rate in Table~\ref{tab:rq2-structural}, SMT validation revealed that these connections were frequently \emph{semantically invalid}, for example by connecting incompatible ports or creating circular dependencies.
This 0\% result shows that large-scale system consistency was not achieved by the evaluated first-pass constrained generation workflow on this benchmark, and therefore marks the point at which structural success no longer translates into system-level correctness. In other words, the multi-file experiment reveals where the remaining difficulty lies: not in producing importable files, but in preserving architectural meaning across files. That distinction motivates the iterative repair workflow supported in ATLAS and points naturally toward stronger design-time guidance, validator-guided revision, and selective human intervention in later stages of the pipeline.

\begin{table}[!htbp]
\centering
\caption{Structural correctness and cross-file reference integrity by complexity tier. All systems achieve perfect file completeness and XSD compliance, while reference resolution degrades with increasing system complexity.}
\label{tab:rq2-structural}

\resizebox{\linewidth}{!}{%
\begin{tabular}{@{}lccccccc@{}}
\toprule
\textbf{Complexity} & \textbf{Systems} & \textbf{Total Files} & \textbf{Completeness} & \textbf{XSD Pass Rate} & \textbf{Total References} & \textbf{Resolved References} & \textbf{Resolution Rate} \\
\midrule
Simple  & 7 & 52  & 52/52 (100\%)  & 52/52 (100\%)  & 605  & 605  & 100.0\% \\
Middle  & 7 & 108 & 108/108 (100\%) & 108/108 (100\%) & 1{,}754 & 1{,}744 & 99.4\% \\
Complex & 6 & 124 & 124/124 (100\%) & 124/124 (100\%) & 2{,}086 & 1{,}936 & 92.8\% \\
\midrule
\textbf{Total} & \textbf{20} & \textbf{284} & \textbf{284/284 (100\%)} & \textbf{284/284 (100\%)} & \textbf{4{,}445} & \textbf{4{,}285} & \textbf{96.4\%} \\
\bottomrule
\end{tabular}%
}
\end{table}

\begin{table}[!htbp]
\centering
\small
\caption{Result analysis findings aggregated by complexity tier. Scores use 5-point scales for quantitative dimensions; C3 uses categorical assessment (minor/moderate/extensive); D1 reports binary import status. Findings illustrate representative failure patterns observed within each tier rather than universal characteristics.}
\label{tab:rq2-analysis}
\begin{tabular}{@{}lp{0.22\textwidth}p{0.58\textwidth}@{}}
\toprule
\textbf{Dimension} & \textbf{Scores} & \textbf{Representative Findings} \\
\midrule
\multicolumn{3}{l}{\textit{Simple Tier (7 systems, 27 components)}} \\
\midrule
Requirements Alignment & A1: 100\%; A2: 100\%; A3: 5/5 & End-to-end functional intent preserved across chain topology. All interfaces emitted as standalone files. Dataflow matches blueprint specifications without omissions. \\
Architectural Quality & B1: 5/5; B2: 5/5 & Linear pipeline topology with correctly paired Provided/Required ports. Sender-Receiver and Client-Server interface kinds consistent with specified types. \\
Engineering Quality & C1: 5/5; C2: 5/5; C3: minor edits & Consistent naming conventions and modular file separation. No architectural changes required. \\
Toolchain Integration & D1: success & All systems import cleanly into commercial AUTOSAR configuration tooling. \\
\midrule
\multicolumn{3}{l}{\textit{Middle Tier (7 systems, 45 components)}} \\
\midrule
Requirements Alignment & A1: 100\%; A2: 100\%; A3: 3/5 & \textit{Representative issue in one system}: Diagnostic coordination module both provides and requires the same Client-Server interface, violating client-server separation. UDS service provided-ports absent on coordination/router components where requirements specified them. \\
Architectural Quality & B1: 2/5; B2: 5/5 & \textit{Representative issue}: Client-Server role confusion in diagnostic subsystem. Local type matching preserved, but global coordination pattern violated. \\
Engineering Quality & C1: 5/5; C2: 5/5; C3: moderate edits & \textit{Representative remedy}: Remove spurious provided services on diagnostic modules; add missing UDS exposure ports. \\
Toolchain Integration & D1: success & All systems import despite semantic inconsistencies. \\
\midrule
\multicolumn{3}{l}{\textit{Complex Tier (6 systems, 54 components)}} \\
\midrule
Requirements Alignment & A1: 100\%; A2: 100\%; A3: 2/5 & \textit{Representative issues}: SecOC chain lacks complete tag generation/verification steps. Health module architecture contradicts observer pattern—pulls business data via required-ports rather than consuming pushed health events. \\
Architectural Quality & B1: 1/5; B2: 5/5 & \textit{Representative issue}: Health module pulls business data via required-ports to infer health, contradicting intended observer pattern where producers push health states via provided-ports. \\
Engineering Quality & C1: 5/5; C2: 5/5; C3: partial rewrite & \textit{Representative remedy}: Refactor health monitoring to pure consumer of health events; reestablish SecOC end-to-end tag flow. \\
Toolchain Integration & D1: success & All systems import successfully. Structural validity does not guarantee semantic correctness. \\
\bottomrule
\end{tabular}
\end{table}

\paragraph{Generation Cost and Scalability.} Blueprint synthesis (Phase 1) incurs moderate cost averaging 17,548 to 26,617 input tokens and 127 to 275 seconds across complexity tiers. Component assembly (Phase 2) dominates computational cost, scaling super-linearly with system complexity: Simple systems average 63,959 output tokens and 195 seconds, Middle systems rise to 184,605 tokens and 530 seconds, and Complex systems reach 280,038 tokens and 717 seconds. The approximate 1.5× token increase and 1.35× time increase from Middle to Complex tiers reflect the combinatorial explosion of cross-file reference formation and context integration requirements as dependency graph density increases.

\paragraph{Result Analysis Findings.} Table~\ref{tab:rq2-analysis} presents aggregated result analysis across the four evaluation dimensions. For Simple systems, all dimensions score maximally. Requirements alignment shows full functional coverage, complete interfaces, and full traceability (A3: 5/5). Architectural quality maintains perfect dependency consistency and type matching (B1/B2: 5/5). Engineering quality remains strong, with consistent naming, modular structure, and only minor edits required (C1/C2: 5/5; C3: minor). All systems also import successfully into commercial tooling (D1: success).

Middle systems maintain perfect functional coverage and interface completeness but exhibit reduced traceability (A3: 3/5). A representative failure pattern observed in one Middle-tier system involves diagnostic coordination components: a diagnostic module incorrectly both provides and requires the same Client-Server interface, violating the client-server separation principle fundamental to AUTOSAR diagnostics architecture. In addition, Unified Diagnostic Services (UDS) exposure is incomplete, with diagnostic provided-ports absent on coordination and router components where the requirements specified them. These issues reduce dependency consistency to B1: 2/5, even though local type matching remains correct at B2: 5/5. Engineering quality stays high, with naming and modularity both at 5/5 and with moderate repair effort required to remove spurious provided services and add missing UDS ports (C3: moderate edits). All Middle systems still import successfully despite these semantic inconsistencies.

Complex systems encounter more severe semantic deviations. A representative Complex-tier system exhibits two major issues: the SecOC (Secure Onboard Communication) authentication chain lacks complete tag generation and verification steps, breaking the end-to-end security guarantee; and health monitoring components contradict the intended observer pattern by pulling business data through required ports instead of consuming health-state events pushed through provided ports. Requirements traceability declines to A3: 2/5, and dependency consistency falls to B1: 1/5, indicating fundamental architectural misalignment. The health-monitoring failure exemplifies architectural role drift: the generated component topology remains locally type-correct at B2: 5/5 but violates the global observer discipline specified in the system requirements. Repair effort escalates to partial rewrite, requiring refactoring of health-monitoring subsystems and rethreading of SecOC tag flows. Despite these issues, all Complex systems import into tooling, demonstrating that structural validity alone does not guarantee semantic correctness for system-level configurations.



\subsubsection{Analysis and Discussion}
\label{sec:rq2-analysis}

\paragraph{What scales cleanly.}
RQ2 shows that the structural part of the pipeline scales substantially better than the semantic part. Across all 20 systems and 284 generated files, ATLAS maintained 100\% file completeness and 100\% XSD compliance. Combined with the strong local type-matching scores in Table~\ref{tab:rq2-analysis}, this indicates that blueprint-guided decomposition plus Layer-1 constrained decoding is effective at preserving file layout, syntactic well-formedness, and many local schema obligations even for comparatively large AUTOSAR assemblies.

\paragraph{What does not scale cleanly.}
The same experiments also show a sharp boundary: structural validity does not imply system-level correctness. Reference resolution remains high (96.4\% overall), yet SHACL and SMT checks expose substantial residual defects, including missing mandatory attributes, role confusion in client/server interactions, broken SecOC chains, and observer-pattern inversions. This is precisely the distinction the paper now tries to make explicit: ATLAS can reliably keep generation inside the \emph{structural} envelope, but global semantic coherence still depends on validators, repair, and in harder cases domain expertise.

\paragraph{Why the blueprint helps, and why it is insufficient.}
The blueprint reduces local search difficulty by committing early to component inventories, filenames, and many reference targets. That explains the strong completeness and importability results. However, the blueprint remains only a partial surrogate for architectural intent. It captures explicit contracts more readily than implicit design disciplines such as service-role separation, end-to-end security threading, or health-monitoring propagation patterns. The middle and complex tiers therefore fail less because the files are malformed than because the local choices remain globally underconstrained. A practical implication is that stronger human intervention at blueprint time, interface freezing for critical contracts, and per-file confirmation at selected decision points could further increase determinism without changing the core ICM/L1/L2 decomposition.

\paragraph{Role of result analysis.}
The findings in Table~\ref{tab:rq2-analysis} are best read as \emph{diagnostic evidence about residual defect burden}. For simple systems, the remaining work is minor. For middle-tier systems, the outstanding issues are often localized and plausibly addressable through validator-guided patches. For complex systems, however, the remaining defects can be architectural, requiring cross-component reinterpretation of intent. The practical implication is that ATLAS usefully shifts the bottleneck away from low-level structural authoring and toward targeted inspection and correction, but it does not by itself establish a human-factors claim about engineering effort reduction once system-level design intent becomes underconstrained.

\paragraph{Implications for the paper's central claim.}
Taken together, RQ1 and RQ2 clarify the level at which ATLAS is currently effective. At single-file scale, the combination of ICM-complete prompting and Layer~1 constrained decoding makes structural correctness highly reliable. At multi-file scale, the same architecture continues to secure completeness, importability, and a strong level of reference organization, but the residual failures shift upward into architectural semantics and cross-file logic. The significance of ATLAS therefore lies in how it reorganizes the problem: structural admissibility becomes routine, while the genuinely difficult part of system construction is surfaced explicitly at the validation and repair boundary. For complex system generation, the natural deployment interpretation is consequently semi-automatic. The ICM provides a shared representation of constraints and evidence, Layer~1 stabilizes structure, and Layer~2 together with targeted human intervention incrementally restores semantic and logical determinism.

% ---------------------------------------------------------------------------------------------------------------------------
%======================================================================
\subsection{Threats to Validity}
\label{sec:threats}
\paragraph{Construct validity.} Our principal outcome measures separate structural validity, semantic/logic validation, and result-analysis-based residual defects. This separation is deliberate, but it also means that no single metric fully captures engineering usefulness. In particular, the result analysis is used as a diagnostic lens on remaining defect burden rather than as a formal productivity study. Because the analysis is not positioned as an external independent assessment, its judgments should be interpreted accordingly.

\paragraph{Internal validity.} The comparison depends on one main base model and a specific serving stack. Although the metamodel-integration + ICM + L2 pipeline is intentionally model-agnostic, quantitative outcomes may shift with different backbones, sampling parameters, or decoding libraries. To reduce this risk, we keep the pipeline fixed across baselines and attribute performance differences to the added grounding and constraint layers rather than to unrelated model changes.

\paragraph{External validity.} AUTOSAR is an appropriate testbed because it combines rich schema information with industrial validation tooling, but it is still only one domain and one particularly schema-rich setting. The current paper therefore supports transfer only by architectural plausibility, not by direct empirical evidence outside AUTOSAR. Likewise, the SHACL/SMT Layer~2 realization used here is one AUTOSAR-specific instantiation of the post-generation validation boundary. Other MDE domains may call for different backends, and mature industrial settings may naturally attach ATLAS to existing validators rather than reconstruct that layer inside the research prototype.


%======================================================================
%======================================================================
\section{Conclusion}
\label{sec:conclusion}
\label{sec:concl}
%======================================================================
ATLAS combines metamodel integration, compiled layered constraints (ICM), and constrained generation plus validation (CVG plus audit-guided repair) to make LLM-based artifact generation more structured, more inspectable, and easier to repair in schema-rich engineering settings. In this paper, that claim is grounded specifically in AUTOSAR rather than in a broad cross-domain validation story.

Across the AUTOSAR evaluation, RQ1 shows that single-file structural reliability is best understood as the combination of ICM-complete prompting and Layer-1 constrained decoding under a fixed, fully instrumented local backbone, while RQ2 shows that the same layered control principles remain useful at multi-file scale in a stronger API-based deployment setting. What does not carry over automatically is full semantic and logical correctness: once generation moves from isolated artifacts to system organization, the remaining difficulty resides in cross-file meaning rather than file construction. The result-analysis findings sharpen this transition. Simple systems often require only limited cleanup, whereas complex systems still demand architectural judgment. The central contribution of ATLAS is thus methodological. It introduces ICM as a unifying constraint representation and uses Layer~2 as a validation boundary through which generated artifacts can be connected to domain-appropriate semantic and logical analysis, including existing tool support when such support is already mature.

\paragraph*{Limitations.}
\begin{itemize}[leftmargin=*,itemsep=0.25ex]
  \item Semantic coverage remains partial. Layer-1 offers strong structural control, but Layer-2 semantic and logical checks still depend on post-generation validation over completed artifacts. In mature domains, the richest realizations of this validation layer may continue to reside in external engineering tools rather than inside the research prototype itself.
  \item Human-factors evidence is still limited. The current result analysis is diagnostic rather than a controlled productivity or usability study, so stronger claims about engineering effort reduction must wait for future work.
  \item Evaluation breadth remains limited. Although the architecture may transfer to other schema-rich settings, this paper evaluates it only in AUTOSAR and does not empirically establish cross-domain effectiveness.
\end{itemize}

\paragraph*{Future Work.}
We see three immediate directions for strengthening the paper's core thesis:
\begin{enumerate}[label=\arabic*),leftmargin=*,itemsep=0.4ex]
  \item Incremental semantic checking during generation, so that at least part of the current post-hoc validation burden can be moved earlier in the pipeline.
  \item Richer repair and adaptation strategies, including audit-guided prompt refinement or fine-tuning that target recurrent high-value failure modes without weakening the explicit constraint layer.
  \item Stronger human-centered evaluation, including controlled studies of repair effort, practitioner confidence, manual-versus-assisted workflows on realistic industrial tasks, and higher-determinism intervention points during blueprinting and per-file generation.
\end{enumerate}

\section*{Data Availability}
The data and materials associated with this study are available at: \url{https://github.com/Abandooon/ATLAS}.

\bibliographystyle{ACM-Reference-Format}
\bibliography{atlas_ref} 

\end{document}



